\documentclass[11pt]{article}
\usepackage[margin=1in]{geometry}
\usepackage{amsmath,amssymb,amsthm,array,booktabs,longtable}
\usepackage[T1]{fontenc}
\usepackage{lmodern}
\usepackage[colorlinks=true,urlcolor=blue,citecolor=blue,linkcolor=blue]{hyperref}
\newtheorem{theorem}{Theorem}
\newtheorem{lemma}[theorem]{Lemma}
\newtheorem{corollary}[theorem]{Corollary}
\newcommand{\C}{\mathbb C}
\newcommand{\T}{\mathbb T}
\newcommand{\Span}{\operatorname{span}_{\C}}
\newcommand{\diag}{\operatorname{diag}}
\title{Ding--Hu, Section V, p.~8 (arXiv:1909.11256v3):\\
negative answers to the two open finite-cover questions}
\author{Haobo Ma\\\small ChronoAI Pte Ltd, Singapore
\and Wenlin Zhang\\\small National University of Singapore\\\small ChronoAI Pte Ltd, Singapore}
\date{November 3, 2026}
\begin{document}
\maketitle
\begin{abstract}
We answer negatively both open questions in Ding and Hu,
\emph{Masking quantum information on hyperdisks}, Section V, p.~8,
arXiv:1909.11256v3 (15 October 2020): the finite-cover questions for regular
subsets of hyperdisks and for nondegenerate masking machines.
An explicit four-coordinate example has span dimension three and requires
continuum many contained hyperdisks, each a single ray.
The same support obstruction gives a complete maskable set requiring
continuum many hyperdisks for every prescribed nondegenerate marginal
spectrum of rank $d\ge4$, at every input dimension $3\le n<d$.
We also determine the cover cardinalities of the two degenerate examples
in the posing paper and distinguish maximal maskable sets of density
matrices from convex hulls of phase states.
\end{abstract}

\section{The question and the example}
Ding and Hu's two open questions in Section V, p.~8 of
arXiv:1909.11256v3 \cite{dinghu} concern finite hyperdisk covers of regular
subsets and of complete maskable sets with nondegenerate marginals.
Here and below a
cover of a set means a union of hyperdisks \emph{contained in that set},
as in their Eq.~(16). Allowing arbitrary supersets would make the question
trivial. A single four-coordinate example answers both questions and
extends to arbitrary positive marginal spectra of rank at least four.

We represent a pure state by its unit vectors; a one-dimensional hyperdisk
is therefore the full circle of representatives of one ray. Passing to
rays gives Ding and Hu's convention. Write $\T=\{z\in\C:|z|=1\}$.
For masking machines we follow their Section III convention: the input
space is restricted to the complex span of its complete maskable set;
$n$ is that dimension and $d$ is the common marginal rank.
Their Definition~1 is the set
\[
 H(\phi,a)=\left\{\sum_{p=1}^m a_pu_p\phi_p:u_p\in\T\right\},
 \qquad a_p>0,\quad \sum_pa_p^2=1,
\]
where the $\phi_p$ are orthonormal. Its span has dimension $m$.
By their Definition~4, $C\subseteq S$ is regular if
$\Span(C)\cap S=C$.

Set
\begin{equation}\label{eq:example}
 \begin{gathered}
 r=\frac{(1,2,3,4)}{\sqrt{30}},\qquad
 S=\{x\in\C^4:|x_j|=r_j\ (0\le j\le3)\},\\
 L=\{x:6x_0+3x_1+2x_2+3x_3=0\},\qquad C=L\cap S.
 \end{gathered}
\end{equation}
The coordinates in $S$ already have squared norm one, and $S$ is a
four-dimensional hyperdisk. The set $C$ is regular: $\Span(C)\subseteq L$
gives $\Span(C)\cap S\subseteq C$, while the reverse inclusion is immediate.

\begin{theorem}\label{thm:main}
Every nonempty hyperdisk contained in $C$ has dimension one. The set $C$
has a continuum of rays and $\Span(C)=L$, of complex dimension three.
Consequently the least cardinality of a hyperdisk cover of $C$ is
$\mathfrak c=|\mathbb R|$.
\end{theorem}

The phase-support argument below is also the mechanism behind the sparse
Gramian description of subhyperdisks in Ding--Hu, Section II, Eq.~(10).
Their Lemma~1 gives finite covers for regular subsets of span dimension
two. The example lies in span dimension three instead. It is also distinct
from their Eq.~(51), which uses partially degenerate marginals. The
no-masking theorem of Modi et al.~\cite{modi} excludes masking every
input state; the present question concerns the structure of the proper
subset that can be masked. The original single-hyperdisk conjecture is
different from the finite-cover question: Ding--Hu's 2019 preprint already
gives a counterexample, and Zhu's 2021 independent argument uses real
states \cite[Section VII, Proposition 4]{zhu}.

\section{The support obstruction}
\begin{lemma}\label{lem:modulus}
If $|\sum_pa_pu_p|$ is constant for all independent $u_p\in\T$, then
at most one $a_p$ is nonzero.
\end{lemma}
\begin{proof}
Fix $p$ and the other phases, and write $B=\sum_{q\ne p}a_qu_q$.
Equality of $|a_pu+B|^2$ at $u=1,-1,i,-i$ forces both real and
imaginary parts of $a_p\overline B$ to vanish. If $a_p\ne0$, then $B=0$
for every choice of the other phases. Comparing all phases equal to one
with the choice $u_q=-1$ forces $2a_q=0$ for each $q\ne p$.
\end{proof}

\begin{proof}[Proof of the hyperdisk assertion in Theorem~\ref{thm:main}]
Let $H(\phi,a)\subseteq C$, and put $v_p=a_p\phi_p$.
Each coordinate of $\sum_pu_pv_p$ has fixed positive modulus $r_j$.
Lemma~\ref{lem:modulus} shows that the nonempty supports
$B_p=\{j:(v_p)_j\ne0\}$ partition $\{0,1,2,3\}$, and
$|(v_p)_j|=r_j$ on $B_p$. The supports are nonempty because $a_p>0$
and $\|\phi_p\|=1$.

Let $\ell(x)=6x_0+3x_1+2x_2+3x_3$.
The identity $\sum_pu_p\ell(v_p)=0$ for all phases, compared with
one phase flipped, gives $\ell(v_p)=0$ for every $p$.
Define
\[
 b_{pj}=\frac{\sqrt{30}}6 c_j(v_p)_j,\qquad (c_0,c_1,c_2,c_3)=(6,3,2,3).
\]
Then $\sum_jb_{pj}=0$, and the nonzero summands have moduli
$1,1,1,2$ at coordinates $0,1,2,3$, respectively.

The block containing coordinate $3$ must contain at least two of
$0,1,2$: one summand of modulus one cannot cancel a summand of modulus
two, and no summands cannot cancel it either. At most one coordinate
remains outside that block. A second nonempty block would consequently
consist of a single nonzero summand and could not have sum zero.
There is exactly one block, so $m=1$, as required.
\end{proof}

\section{Distinct rays and the span}
For $-3/4\le t\le-1/2$, define
\begin{equation}\label{eq:phases}
 \begin{gathered}
 z=t+i\sqrt{1-t^2},\quad w=-2-z,\quad
 q=5+4t,\quad k=\sqrt{\frac{4-q}{q}},\\
 u(t)=\left(z,\frac{w(1+ik)}2,\frac{w(1-ik)}2,1\right).
 \end{gathered}
\end{equation}
Here $2\le q\le3$, so every square root is real and the denominator is
positive. We have $|z|^2=1$, $|w|^2=q$, and
\[
 \left|\frac{w(1\pm ik)}2\right|^2
 =\frac{q(1+k^2)}4=1,\qquad u_0(t)+u_1(t)+u_2(t)+2u_3(t)=0.
\]
Thus $x(t)=(r_ju_j(t))_{j=0}^3\in C$.
If $x(t)=\alpha x(t')$, the fixed nonzero fourth coordinate forces
$\alpha=1$. The real part of the first coordinate then forces $t=t'$.
The interval therefore injects into the rays of $C$. There are at most
continuum many such rays because $C\subseteq\C^4$, so there are exactly
continuum many. Each hyperdisk contained in $C$ covers one ray, while
all singleton rays are hyperdisks. This proves the exact cover cardinality.

For completeness, the span is exactly the displayed hyperplane.
Put $\omega=(-1+i\sqrt3)/2$ and let $D=\diag(r_0,r_1,r_2,r_3)$.
The three vectors
\begin{equation}\label{eq:basis}
 D(-1,\omega,\overline\omega,1),\qquad
 D(-1,\overline\omega,\omega,1),\qquad
 D(\omega,-1,\overline\omega,1)
\end{equation}
belong to $C$ since $\omega+\overline\omega=-1$.
In a linear dependence with coefficients $\alpha,\beta,\gamma$,
the fourth coordinate gives $\alpha+\beta+\gamma=0$ and the first gives
$(1+\omega)\gamma=0$. Hence $\gamma=0$ and $\beta=-\alpha$.
The second coordinate then gives $(\omega-\overline\omega)\alpha=0$,
so all coefficients vanish. These vectors are independent, and
$\dim_{\C}L=3$, completing the proof of Theorem~\ref{thm:main}.

\begin{corollary}\label{cor:old-bound}
The proposed span-dimension upper bound in Ding--Hu,
arXiv:1909.11256v1, Section V, p.~6 (25 September 2019), is false:
a regular subset of span dimension three can require $\mathfrak c$
contained hyperdisks.
\end{corollary}
\begin{proof}
Apply Theorem~\ref{thm:main} to $C$ in Eq.~\eqref{eq:example}.
\end{proof}

\begin{corollary}\label{cor:minimal}
Four is the smallest ambient hyperdisk dimension admitting a regular
subset without a finite cover by contained hyperdisks.
For nondegenerate masking machines, the smallest input dimension and
marginal rank allowing such a cover obstruction are $n=3$ and $d=4$.
\end{corollary}
\begin{proof}
For an ambient hyperdisk of dimension at most three, a regular subset
of the same span dimension is the whole hyperdisk, by Ding--Hu Eq.~(15).
One of span dimension two is a hyperdisk or two rays, by their Lemma~1.
A nonempty subset of span dimension one is a single ray, since regularity
includes all its unit representatives. A subset of span dimension zero
is empty and has the empty cover. These cases exhaust the possibilities.
For nondegenerate masking, the complete target set is a regular subset
of a $d$-dimensional Schmidt hyperdisk by Ding--Hu Eqs.~(26)--(29).
Lemma~1 also excludes any obstruction of span dimension at most two,
regardless of $d$. Corollary~\ref{cor:masker} attains $(n,d)=(3,4)$.
\end{proof}

\section{A nondegenerate masking machine}
\begin{corollary}\label{cor:masker}
There is a masking machine with input dimension $n=3$ and Schmidt
number $d=4$, whose complete maskable set has no finite or countable
hyperdisk cover, although its fixed marginal states are nondegenerate.
\end{corollary}
\begin{proof}
Take the input Hilbert space to be $L$ with its inherited inner product,
and define the isometry
\[
 V:L\longrightarrow\C^4\otimes\C^4,\qquad
 Vx=\sum_{j=0}^3x_j|j\rangle|j\rangle.
\]
For every unit $x\in L$, either partial trace of $|Vx\rangle\langle Vx|$
is $\diag(|x_0|^2,|x_1|^2,|x_2|^2,|x_3|^2)$; the off-diagonal
terms disappear because $\langle j|k\rangle=0$ for $j\ne k$.
Both marginals equal
\[
 \rho=\frac1{30}\diag(1,4,9,16)
\]
if and only if $x\in L\cap S=C$. Thus $C$ is the \emph{complete}
maskable input set, satisfying both clauses of Ding--Hu Definition~5,
not merely a selected subset. The eigenvalues of $\rho$ are positive
and pairwise distinct, so $d=4$ and the marginals are nondegenerate.
By Theorem~\ref{thm:main}, the input dimension is three. An isometry
preserves inner products and spans, hence takes hyperdisks to hyperdisks
and does so in reverse on its image. The same cover obstruction holds
for the target set $V(C)$. This proves the negative answer to the first
question of Ding--Hu Section V.
\end{proof}

\begin{corollary}\label{cor:spectra}
Let $d\ge4$ and $r_j>0$ with $\sum_{j=0}^{d-1}r_j^2=1$. Set
\[
 S_d=\{x\in\C^d:|x_j|=r_j\},\qquad
 L_d=\left\{x:\frac{x_0}{r_0}+\frac{x_1}{r_1}
                 +\frac{x_2}{r_2}+\frac{2x_3}{r_3}=0\right\}.
\]
Then $C_d=L_d\cap S_d$ is regular, spans $L_d$, and its least
hyperdisk-cover cardinality is $\mathfrak c$.
The isometry $x\mapsto\sum_jx_j|jj\rangle$ on $L_d$ has complete maskable
set $C_d$ for marginals $\diag(r_0^2,\ldots,r_{d-1}^2)$ and input dimension
$d-1$. In particular, every nondegenerate rank-$d$ spectrum occurs.
More generally, for every $3\le n<d$, the subspace
\[
 L_{d,n}=\left\{x\in L_d:
       \frac{x_j}{r_j}=\frac{x_3}{r_3}\quad(n+1\le j\le d-1)\right\}
\]
has dimension $n$, and $L_{d,n}\cap S_d$ has the same cover cardinality
and complete masking realization. Hence the pairs $(n,d)$ admitting a
nondegenerate masking machine without a finite hyperdisk cover are
exactly $3\le n<d$.
\end{corollary}
\begin{proof}
Regularity follows as in Eq.~\eqref{eq:example}. For $d=4$, diagonal
rescaling of the phase vectors in Eqs.~\eqref{eq:phases} and
\eqref{eq:basis} proves the continuum and span assertions for any $r$.
For $d>4$, append fixed phases in the remaining coordinates. Flipping
one of these phases and subtracting the two resulting vectors gives
a nonzero multiple of its coordinate vector. Together with the three
rescaled vectors in Eq.~\eqref{eq:basis}, these give $d-1$ independent
directions in $L_d$.

For a hyperdisk $H\subseteq C_d$, Lemma~\ref{lem:modulus} again assigns
each coordinate to exactly one phase generator. Every generator satisfies
the defining linear equation of $L_d$. Among the first four coordinates,
the block containing coordinate $3$ needs at least two unit-weight
coordinates; any remaining coordinate cannot form another zero-sum block.
Thus all four coordinates belong to one generator. Projection of $H$
onto the first four coordinates is therefore contained in a single ray.
The rescaled family in Eq.~\eqref{eq:phases}, with fixed appended phases,
has continuum many different projected rays. Each contained hyperdisk
covers at most one member of this family, so any cover has cardinality
at least $\mathfrak c$. Singleton-ray hyperdisks give the reverse bound
because $C_d\subseteq\C^d$. The marginal calculation in
Corollary~\ref{cor:masker} applies without change.

For $L_{d,n}$, keep coordinates $4,\ldots,n$ free and tie every later
coordinate to coordinate $3$ as displayed. The three phase vectors and
the sign-flip differences in the $n-3$ free coordinates span an
$n$-dimensional subspace; the constraints give the matching upper bound.
The continuum family still satisfies the ties. Every contained hyperdisk
is contained in $C_d$, so the projected-ray obstruction and cover
cardinality persist. Finally, Ding--Hu Lemma~1 excludes $n\le2$,
Eqs.~(26)--(29) give $n\le d$ in the nondegenerate case, and equality
$n=d$ gives the whole Schmidt hyperdisk. These prove the necessity.
\end{proof}

\section{Consequences for the degenerate examples}
The nondegenerate construction shows that degeneracy is not necessary
for an infinite least cover cardinality. Ding--Hu already exhibited
degenerate examples in Eqs.~(36) and (51); their precise cover
cardinalities follow as well.

\begin{corollary}\label{cor:degenerate}
Each complete target set in Ding--Hu Eqs.~(36) and (51) has least
hyperdisk-cover cardinality $\mathfrak c$. The latter has no contained
hyperdisk of dimension greater than one, as stated in that paper.
\end{corollary}
\begin{proof}
For Eq.~(51), take $\lambda_1,\lambda_2>0$ with
$2\lambda_1+\lambda_2=1$ and $\lambda_1\ne\lambda_2$.
In the orthonormal basis
\[
 |00\rangle,\quad |11\rangle,\quad
 \frac{i\sqrt{\lambda_1/2}(|01\rangle+|10\rangle)
             +\sqrt{\lambda_2}|22\rangle}{\sqrt{\lambda_1+\lambda_2}},
\]
its representatives have coefficients
$(ae^{i\eta},ae^{-i\eta},b)$, up to a common phase, where
$a=\sqrt{\lambda_1/2}$ and $b=\sqrt{\lambda_1+\lambda_2}$.
They satisfy $x_0x_1=(a^2/b^2)x_2^2$ and have fixed positive moduli.
For any contained hyperdisk, the disjoint-support lemma assigns
coordinates $0,1,2$ to generators $p_0,p_1,p_2$.
Independent phase variation in this quadratic identity forces
$u_{p_0}u_{p_1}=u_{p_2}^2$ for all phases, hence
$p_0=p_1=p_2$. There is one generator. The fixed nonzero third
coefficient distinguishes continuum many rays as $\eta$ ranges over
$[0,2\pi)$, proving the assertion for Eq.~(51).

For Eq.~(36), use its target span
$W=\Span\{|00\rangle,|11\rangle,|01\rangle+|10\rangle\}$.
Every hyperdisk contained in its maximally entangled target set has
dimension at most two. Indeed, write its weighted orthonormal generators
as $2\times2$ coefficient matrices $A_p$, with $*$ denoting conjugate
transpose. Constant marginals for all
independent phases force $A_pA_q^*=A_p^*A_q=0$ for $p\ne q$, by the
coefficients of $u_p\overline{u_q}$. Thus the nonzero ranges of $A_p$
are pairwise orthogonal in $\C^2$.
The orthonormal vectors
\[
 f_0=\frac{i(|00\rangle+|11\rangle)}{\sqrt2},\quad
 f_1=\frac{|00\rangle-|11\rangle}{\sqrt2},\quad
 f_2=\frac{|01\rangle+|10\rangle}{\sqrt2}
\]
give a continuum of target rays
$[f_0+t f_1+t^2 f_2]$ for $t\in\mathbb R$: normalization divides by
$\sqrt{1+t^2+t^4}$, and either marginal is $I/2$.
Any complex subspace of $W$ of dimension at most two is annihilated
by a nonzero linear functional. On this family its value is a nonzero
polynomial of degree at most two, so it vanishes at at most two parameters.
Every contained hyperdisk covers at most two of these rays.
A family of fewer than $\mathfrak c$ such finite intersections cannot
cover $\mathbb R$; singleton rays again give the matching upper bound.
\end{proof}

\section{A related maximality question}
Zhu's Section VII, Conjecture~1 \cite{zhu}, concerns the \emph{convex hull}
of the pure phase states, not the finite-cover question. Its distinction
from the full fixed-diagonal set matters in dimension four and above.
The fixed-diagonal set was already characterized as the complete maskable
set of the diagonal isometry in the April 2020 preprint of
Du et al.\ \cite[arXiv version, Theorem 3.2]{du}.
The convex-geometric ingredients are the classical correlation-matrix
rank bound and construction of Li and Tam
\cite[Theorem 1, Corollary 2 and Section 2.1]{litam}; the proof below
specializes them explicitly.
For a positive normalized vector $r$, write
\[
 \mathcal D_r=\{\sigma\ge0:\diag(\sigma)=(r_j^2)_j\},\qquad
 \mathcal P_r=\{|x\rangle\langle x|:|x_j|=r_j\}.
\]
\begin{corollary}\label{cor:maximal}
The set $\mathcal D_r$ is maximal maskable among density matrices on
$\C^d$. The convex hull of $\mathcal P_r$ equals $\mathcal D_r$ for
$d\le3$ and is a proper subset for every $d\ge4$.
Consequently Zhu's Conjecture~1 holds for $d=2,3$ and is false for $d\ge4$;
the corrected maximal set is $\mathcal D_r$.
\end{corollary}
\begin{proof}
The isometry $V:\C^d\to\C^d\otimes\C^d$ defined on the full input
space by $V|j\rangle=|jj\rangle$ masks $\mathcal D_r$.
For any isometry masking $\mathcal P_r$, let $A_j$
be the coefficient matrix of the image of $|j\rangle$.
Fourier coefficients of either constant marginal give
$A_jA_k^*=A_j^*A_k=0$ for $j\ne k$, since $r_jr_k>0$.
Thus $B_j=A_jA_j^*$ are nonzero, pairwise orthogonal positive matrices,
and the first marginal of an arbitrary input $\sigma$ is
$\sum_j\sigma_{jj}B_j$. It equals $\sum_jr_j^2B_j$ only if
$\sigma_{jj}=r_j^2$ for every $j$, as multiplication by $B_j$ and
taking traces shows. Hence no masker of $\mathcal P_r$, and therefore
none of $\mathcal D_r$, can include a state outside $\mathcal D_r$.

Write $D=\diag(r)$ and $\sigma=DGD$; then $G\ge0$ has diagonal one.
If $G=UU^*$ has rank $s\ge2$ and $d\le3$, the $s^2$-dimensional real
space of Hermitian $s\times s$ matrices contains a nonzero $Q$ with
$u_jQu_j^*=0$ for every row $u_j$ of $U$, since $s^2>d$.
The rows span $\C^s$, so $Q$ is neither positive nor negative
semidefinite. Moving $I+\epsilon Q$ in both directions to its two
positive-semidefinite endpoints decomposes $G$ into a convex combination
of two smaller-rank matrices with diagonal one. Induction reduces to
rank-one phase matrices. This proves the equality for $d\le3$.

For $d\ge4$, let the first four rows of $U$ be
\[
 (1,0),\quad(0,1),\quad(1,1)/\sqrt2,\quad(1,i)/\sqrt2,
\]
and let every further row be $(1,0)$. Then $\sigma=DUU^*D\in\mathcal D_r$.
If $\sigma$ were a convex combination of elements
$|x\rangle\langle x|\in\mathcal P_r$, positivity would put every such
$x$ in $\operatorname{range}(DU)$. Write $x=DU(a,b)^T$.
The first two coordinate moduli give $|a|=|b|=1$; the next two give
$|a+b|^2=|a+ib|^2=2$. Hence both real and imaginary parts of
$a\overline b$ vanish, contradicting $|a\overline b|=1$.
Thus $\sigma$ lies outside that convex hull, which is consequently
not maximal maskable.
\end{proof}

\section{Coverage and verification}
The coverage uses the source's span-dimension convention for hyperdisks.
The two open questions are refuted; the other entries below identify
proved consequences or the earlier results used for comparison.
{\small
\begin{longtable}{@{}p{.34\linewidth}p{.60\linewidth}@{}}
\toprule Source item & Status and conclusion \\\midrule\endhead
Ding--Hu Section V, p.~8, second open question; Eq.~(16) &
Refuted-and-corrected: Theorem~\ref{thm:main} gives optimal cardinality
$\mathfrak c$, not a finite cover.\\
Section V, first open question &
Refuted-and-corrected: Corollaries~\ref{cor:masker} and \ref{cor:spectra}
give nondegenerate examples for every rank $d\ge4$ and prescribed spectrum.\\
arXiv:1909.11256v1, Section V, p.~6, proposed dimension bound &
Refuted-and-corrected: Corollary~\ref{cor:old-bound} gives span dimension
three and least cover cardinality $\mathfrak c$.\\
The same v1 conclusion, question on mixed-state structure &
Not covered as a general classification. Du et al.\ \cite{du}
characterize the diagonal-isometry case; Corollary~\ref{cor:maximal}
determines its maximality among density matrices.\\
Section II, Eq.~(10), Lemma~1 and Eq.~(15) &
Prior proof cited: disjoint supports and the low-dimensional classification.
Proved consequence: sharp dimensions in Corollary~\ref{cor:minimal}.\\
Section III A, Eqs.~(30)--(32); original single-hyperdisk conjecture &
Prior proof cited: Ding--Hu's counterexample, already in the 2019 preprint;
Zhu's 2021 version supplies a later independent real-state argument.\\
Section III B, Eq.~(36); Section IV, Eq.~(51) &
Prior proof cited for the examples; proved exact cover cardinalities
in Corollary~\ref{cor:degenerate}.\\
Section IV, Theorem~1; Theorem~2 and Eq.~(50) &
Prior proof cited: qubit classification, conditional $n=d=3$ classification,
and three-ray example. Their statements are unchanged.\\
Zhu Section VII, Conjecture~1, arXiv:2010.07843v2, p.~9 &
Refuted-and-corrected for $d\ge4$, proved for $d=2,3$:
Corollary~\ref{cor:maximal} distinguishes convex phase hulls and
fixed-diagonal maximal sets.\\
Du et al., arXiv:2004.14572v1, Theorem 3.2; Li--Tam, Theorem 1,
Corollary 2 and Section 2.1 &
Prior proof cited: fixed-diagonal maskability (2020), correlation-matrix
rank bound and explicit rank-two construction (1994).\\
\bottomrule
\end{longtable}}

The pinned Lean source is
\href{https://github.com/the-omega-institute/trureturing/blob/01b573fd9b691a4be3a8b0ecbb41b3311eda19b5/D5/S3/Quantum/Masking/HyperdiskRegularSubsetCover.lean}{\texttt{HyperdiskRegularSubsetCover.lean}}.
Its declarations lie in the namespace
\begin{center}\small
\texttt{D5.S3.Quantum.Masking.HyperdiskRegularSubsetCover}.
\end{center}
Here \texttt{IsHyperdisk} and \texttt{IsRegularSubset} encode the definitions
above, \texttt{claim} is the universal finite-cover assertion, and the
public declaration \texttt{result} proves \(\neg\texttt{claim}\).
The file's auxiliary \texttt{phases\_unit} establishes phase membership
throughout the displayed interval; \texttt{family\_ray\_injective} separates
the countable family at $t_n=-1/2-1/(4(n+1))$; and
\texttt{hyperdisk\_single\_ray} places any two members of a contained
hyperdisk on the same ray. These are local auxiliary lemmas in the pinned
file.

The continuum cardinality, explicit span equality, minimal dimensions,
all masking realizations, degenerate comparisons and maximality assertions
have written proofs here. The accompanying \texttt{verify.py} checks
all 15 four-coordinate support partitions, exact phase identities at
32 rational parameters, a nonzero span minor, finite arbitrary-spectrum
examples, and the four-row rank-two matrix used in
Corollary~\ref{cor:maximal}. The finite checks complement, rather than
replace, the arguments for all parameters and dimensions.

\paragraph{License.} This note is licensed under Creative Commons
Attribution 4.0 International (CC BY 4.0).

\begin{thebibliography}{9}
\bibitem{dinghu}
F. Ding and X. Hu, \emph{Masking quantum information on hyperdisks},
Phys. Rev. A \textbf{102}, 042404 (2020),
\href{https://doi.org/10.1103/PhysRevA.102.042404}{doi:10.1103/PhysRevA.102.042404}.
\href{https://arxiv.org/abs/1909.11256v3}{arXiv:1909.11256v3},
Definitions 1, 4 and 5, Lemma 1, and Section V, p.~8.
\bibitem{modi}
K. Modi, A. K. Pati, A. Sen(De), and U. Sen,
\emph{Masking quantum information is impossible},
Phys. Rev. Lett. \textbf{120}, 230501 (2018),
\href{https://doi.org/10.1103/PhysRevLett.120.230501}{doi:10.1103/PhysRevLett.120.230501}.
\bibitem{zhu}
H. Zhu, \emph{Hiding and masking quantum information in complex and real
quantum mechanics}, Phys. Rev. Research \textbf{3}, 033176 (2021),
\href{https://doi.org/10.1103/PhysRevResearch.3.033176}{doi:10.1103/PhysRevResearch.3.033176}.
\href{https://arxiv.org/abs/2010.07843v2}{arXiv:2010.07843v2},
Section VII, p.~9, Conjecture~1 and Propositions 3--4.
\bibitem{du}
Y. Du, Z. Guo, H. Cao, K. Han and C. Yang,
\emph{Masking quantum information encoded in pure and mixed states},
Int. J. Theor. Phys. \textbf{60}, 2380--2399 (2021; online 26 July 2020),
\href{https://doi.org/10.1007/s10773-020-04542-w}{doi:10.1007/s10773-020-04542-w}.
\href{https://arxiv.org/abs/2004.14572v1}{arXiv:2004.14572v1}
(30 April 2020), Theorem 3.2; the preprint lists Cao first.
\bibitem{litam}
C.-K. Li and B.-S. Tam, \emph{A note on extreme correlation matrices},
SIAM J. Matrix Anal. Appl. \textbf{15}, 903--908 (1994),
\href{https://doi.org/10.1137/S0895479892240683}{doi:10.1137/S0895479892240683}.
\href{https://cklixx.people.wm.edu/corr.pdf}{Author-hosted text},
Theorem 1, Corollary 2 and Section 2.1, pp.~904--905.
\end{thebibliography}
\end{document}
